\subsection{Completed Input and Selection Interventions}
\label{sec:kitti-interventions}
The completed interventions changed input capacity, patch support,
generated-point dose, and selection rules. Their separate baselines
and evaluation scopes are retained below. Differences between tables
are not interpreted as isolated method effects when the pipelines
differ.

\subsubsection{Removing Entrance Resampling}
The all-valid PointRCNN experiment evaluated the original cloud and
four generated Line~A inputs without standard entrance resampling.
Table~\ref{tab:k-all-valid} shows that the original-input Moderate
result itself fell from about 81.95 to 74.91. PDANS was the strongest
generated input at 63.49, but still trailed the all-valid original.
Removing a point cap changed the representation expected by the
detector rather than simply releasing a performance bottleneck.

\begin{table}[!t]\centering\setlength{\tabcolsep}{4pt}
\begin{tabular}{lrrr}\toprule
Input & Easy & Moderate & Hard\\\midrule
Original / standard & 92.26 & 81.95 & 77.75\\
Original / all-valid & 79.83 & 74.91 & 72.73\\
PDANS / all-valid & 83.59 & 63.49 & 54.34\\
PU-GCN / all-valid & 78.21 & 54.26 & 47.11\\
PU-EF / all-valid & 67.11 & 43.85 & 37.10\\
PU-Net old / all-valid & 12.23 & 8.82 & 8.17\\
\bottomrule\end{tabular}
\caption{Recorded full-validation all-valid-input experiment. Values are Car 3D AP from that historical protocol. PU-EF is the compatibility path and PU-Net old retains the faulty adapter.}\label{tab:k-all-valid}\end{table}

\subsubsection{Correcting Local Patch Support across Methods}
Table~\ref{tab:k-patch-pairs} retains all four paired surface-patch
pilots and their two batch baselines. Every method improved with
surface-c32. Fixed-normalisation PU-Net additionally reached 44.0985
with cover-kNN, against its batch baseline of 81.7441. Local support
therefore affected the recorded performance, but did not close the
baseline gap.

\begin{table}[!t]\centering\setlength{\tabcolsep}{4pt}
\begin{tabular}{llrrr}\toprule
Batch & Method & Baseline & Old patch & Surface\\\midrule
Batch 1 & PDANS & 81.56 & 57.99 & 68.47\\
Batch 1 & PU-GCN & 81.56 & 55.81 & 63.45\\
Batch 2 & PU-EF & 81.74 & 41.74 & 58.02\\
Batch 2 & PU-Net fixed & 81.74 & 16.81 & 22.93\\
\bottomrule\end{tabular}
\caption{PointRCNN 256-frame paired patch pilots, Car Moderate 3D AP. Surface denotes surface-c32; the two pilot batches retain separate baselines.}\label{tab:k-patch-pairs}\end{table}

\subsubsection{Region Splitting with Complete Core-Point Retention}
The region-split pilot retained every core point across 2,067 shared
regions with a 3~m halo and at most 16,384 points per detector input.
Table~\ref{tab:k-region-all} reports every method and class.
PDANS was closest to the baseline, but no generated input improved
the three principal class outcomes together.

\begin{table}[!t]\centering\setlength{\tabcolsep}{4pt}
\begin{tabular}{lrrr}\toprule
Input & Car & Pedestrian & Cyclist\\\midrule
Baseline & 77.82 & 58.21 & 78.29\\
PDANS & 68.87 & 58.07 & 71.58\\
PU-GCN & 61.51 & 55.47 & 57.48\\
PU-EF & 61.28 & 56.14 & 56.18\\
PU-Net fixed & 26.18 & 25.02 & 28.92\\
\bottomrule\end{tabular}
\caption{Region-split pilot on 256 frames: Moderate 3D AP$_{R40}$. Generated rows use surface-c32; PU-Net normalisation is fixed. This is a separate three-class PointRCNN implementation and inference protocol.}\label{tab:k-region-all}\end{table}

An independent 64-frame region-split V4 holdout improved
Car/Pedestrian/Cyclist from 75.7267/38.1070/15.0000 to
77.4777/45.8565/17.2222. These positive results remain part of the
completed evidence, bounded by the small holdout and its region-split
protocol. In native whole-frame OpenPCDet PointRCNN, V4 changed Car
by $+4.0651$ under deterministic workers=0 but $-2.2718$ under
workers=2, with other class effects also changing sign. The records
therefore do not establish improvement stable across those sampling
implementations.

\subsubsection{Dose and Selector Components}
The 256-frame PDANS Line~A dose sequence used 2.5\%, 5\%, 7.5\%,
and 10\% generated points. Its reported Moderate results were
approximately 85.10, 82.37, 82.25, and 82.42, against 82.66.
All tested Line~B doses remained below the sparse reference.
Full-validation g10 runs completed for all four methods:
Line~A PDANS/PU-GCN/PU-EF/PU-Net old reached approximately
79.951/79.478/76.433/72.371, and Line~B reached
59.048/55.954/52.415/45.700. The g25 and g50 runs did not complete
and are not treated as evaluated results.

The V1 component experiment compared density, fixed dose, matched-real,
nearest-neighbour, and disabled-component variants
(Table~\ref{tab:k-components}). The full variant had the highest
single-evaluation score, but the matched-real control also improved.
The later evaluator-seed probe described above shows why this one
favourable score did not establish a stable gain.

\begin{table}[!t]\centering\setlength{\tabcolsep}{4pt}
\begin{tabular}{lr}\toprule
Configuration & Moderate\\\midrule
Baseline & 81.18\\
Density-only & 79.86\\
Fixed 2.5\% PDANS & 81.44\\
Full & 82.40\\
Matched-real dynamic & 81.80\\
NN-only & 81.59\\
No-adaptive & 80.15\\
No-confidence & 81.64\\
No-sparse & 80.94\\
\bottomrule\end{tabular}
\caption{PointRCNN V1 selector/component pilot on 256 frames, Car Moderate 3D AP$_{R40}$. The separate 16-seed probe qualifies these single-evaluation scores.}\label{tab:k-components}\end{table}

\subsubsection{Detector-Aware Selection Versions}
V2 used final-box guidance, V3 used internal pre-RCNN or pre-NMS
proposals, V4 anchored measured voxels, and V5 introduced class-adaptive
gating. PointRCNN V2 voxel-only and proposal-only scores were
79.3805 and 75.7533; their combined score was 79.2269.
Internal proposals in V3 yielded 79.2933, so that replacement did
not remove the loss. Table~\ref{tab:k-versions} retains the detector
and class comparisons.

\begin{table}[!t]\centering\setlength{\tabcolsep}{4pt}
\begin{tabular}{lrrrr}\toprule
Version & PR Car & CP Car & CP Ped. & CP Cyc.\\\midrule
Baseline & 81.18 & 79.84 & 43.51 & 76.11\\
V2 full & 79.23 & 76.77 & 48.56 & 72.28\\
V3 full & 79.29 & 76.74 & 48.59 & 72.29\\
V4 & 81.49 & 79.82 & 43.07 & 75.77\\
\bottomrule\end{tabular}
\caption{Completed 256-frame PDANS versions: PointRCNN Car and CenterPoint Car/Pedestrian/Cyclist Moderate 3D AP. CenterPoint Pedestrian/Cyclist baselines are obtained from the recorded scores and baseline deltas.}\label{tab:k-versions}\end{table}

V4 reduced the median number of newly activated voxels to zero,
yet did not improve all classes. V5 changed development-set
CenterPoint Car/Pedestrian/Cyclist by
$-0.0184/+0.0004/+0.0000$ and reused the PointRCNN V4 result.
The completed independent 64-frame holdout did not retain a positive
Car effect: 74.1881 became 74.0512, while BEV fell from 86.3893
to 85.2910. Input control recovered performance during development;
the holdout limits the claim of a general benefit.

\subsubsection{Line B Selection and Small Negative Probes}
Table~\ref{tab:k-lineb-selectors} compares the completed c2048-r4,
E1, consensus, and random pilots. For PointRCNN, PDANS favoured E1,
whereas PU-GCN and fixed PU-Net had their highest listed score with
consensus. For CenterPoint, consensus reduced PDANS, PU-GCN, and
PU-EF relative to random selection but improved PU-Net.
No listed generated input reached its corresponding baseline.

\begin{table}[!t]\centering\setlength{\tabcolsep}{4pt}
\begin{tabular}{lrrrrr}\toprule
Protocol & Base & PDANS & PU-GCN & PU-EF & PU-Net\\\midrule
\shortstack[l]{PointRCNN\\c2048-r4} & 64.04 & 36.62 & 24.96 & 21.24 & 2.40\\
\shortstack[l]{PointRCNN\\E1} & 64.04 & 44.05 & 32.84 & 28.26 & 10.58\\
\shortstack[l]{PointRCNN\\consensus} & 63.42 & 39.25 & 33.16 & 22.44 & 18.88\\
\shortstack[l]{CenterPoint\\random} & 61.66 & 47.74 & 41.38 & 35.43 & 15.15\\
\shortstack[l]{CenterPoint\\consensus} & 61.66 & 43.61 & 37.31 & 27.33 & 19.80\\
\bottomrule\end{tabular}
\caption{Line~B 256-frame selection pilots, Car Moderate 3D AP. PU-EF is the compatibility implementation and PU-Net uses fixed normalisation. Protocol-specific baselines are not pooled.}\label{tab:k-lineb-selectors}\end{table}

A quality-selection probe used three frames without a voxel quota
and five frames with quota 6. Without a quota, reported quality
precision was 0.6821 versus 0.8155 for random selection, and unsupported
fraction was 0.5825 versus 0.4940. With quota 6, the comparisons
were 0.5965 versus 0.8024 and 0.6329 versus 0.5337.
These small geometric probes had no detector AP and were not expanded.
Their negative outcomes qualify the selection hypothesis, while saying
nothing directly about detector performance.

Together, these interventions tested capacity, local support, dose,
proposal guidance, measured-voxel anchoring, and selector components.
Several recovered performance within their own protocols, but no
consistent cross-protocol gain survived the available controls.
They motivated the later PU-GCN input construction and detector
adaptation, whose final full-validation results remain separate.
